% J20 — Mathieu M_22 Substrate-Prime: Order-Factorization Coincidences
%
% Brayden R. Sanders, M. Gish
% Submitted to: American Mathematical Monthly
% Verification: m22_decomposition.py (PASS at machine precision)
% Companions: foundational; cited by later J's
% MSC 2020: 20D08, 05B05, 11A41, 20B25
%
% NOTE: A separate manuscript file WP_PARADOX_CLASSIFIER.md preserved in
%   this folder is content for J53 (Phase 5 paradox classifier
%   expository); the present M_22 substrate-prime manuscript is what
%   this J20 slot ships.

\documentclass[11pt,reqno]{amsart}

\usepackage{amsmath, amssymb, amsthm, mathtools}
\usepackage[margin=1in]{geometry}
\usepackage[colorlinks=true,linkcolor=blue,citecolor=blue,urlcolor=blue]{hyperref}
\usepackage{microtype}
\usepackage{booktabs}
\usepackage{tabularx}

\theoremstyle{plain}
\newtheorem{theorem}{Theorem}[section]
\newtheorem{proposition}[theorem]{Proposition}
\newtheorem{lemma}[theorem]{Lemma}
\newtheorem{corollary}[theorem]{Corollary}

\theoremstyle{definition}
\newtheorem{definition}[theorem]{Definition}

\theoremstyle{remark}
\newtheorem*{remark}{Remark}

\newcommand{\Z}{\mathbb{Z}}
\newcommand{\Mtwentytwo}{M_{22}}
\newcommand{\sig}{\sigma}

\title[Mathieu $\Mtwentytwo$ Substrate-Prime]
{Mathieu $\Mtwentytwo$ Substrate-Prime: \\
 Order-Factorization Coincidences}

\author{Brayden R.\ Sanders}
\address{7Site LLC, Hot Springs, Arkansas, USA}
\email{brayden@7site.co}

\author{M.\ Gish}
\address{Independent Researcher, Hot Springs, Arkansas, USA}
\email{monica.gish1992@gmail.com}

\subjclass[2020]{20D08, 05B05, 11A41, 20B25}
\keywords{Mathieu group, sporadic simple group, Steiner system,
substrate prime, order factorization, hexad}

\date{\today}

\begin{document}

\begin{abstract}
We document a quantitative non-genericity phenomenon between the
sporadic Mathieu group $\Mtwentytwo$ (order $|\Mtwentytwo| = 443\,520
= 2^{7} \cdot 3^{2} \cdot 5 \cdot 7 \cdot 11$) and a discrete
substrate $(\Z/10\Z, \sig, W)$ defined in \cite{SandersGishFourCore}.
The substrate distinguishes the prime set $\{2, 3, 5, 7, 11\}$ ---
the same primes (with the same multiplicities) appearing in
$|\Mtwentytwo|$.

\textbf{Main observation (Theorem~\ref{thm:nongenericity}).} Of
$\Mtwentytwo$'s $12$ irreducible complex representations, exactly
$10$ have dimensions whose only prime factors lie in
$\{2, 3, 5, 7, 11\}$ with at most one factor of $2$; equivalently,
\emph{seven} non-trivial irrep dimensions ($21, 45, 45, 55, 99, 231,
385$) factor strictly in $\{3, 5, 7, 11\}$, and two more ($154 = 2
\cdot 7 \cdot 11$ and $210 = 2 \cdot 3 \cdot 5 \cdot 7$) admit a
single factor of $2$. The fraction $10/12 \approx 83\%$ contrasts
with the null-model density $67/385 \approx 17.4\%$ of integers in
$[1, 385]$ that satisfy the same factorization condition; under a
naive null where the $12$ dimensions were uniformly drawn from
$[1, 385]$, the observed concentration occurs with probability
$\approx 1.2 \times 10^{-6}$.

A short table presents the standard parameters of the Steiner system
$S(3, 6, 22)$ (block count $b = 77 = 7 \cdot 11$, replication
$r = 21 = 3 \cdot 7$, pair-replication $\lambda_{2} = 5$, block size
$k = 6$) with their substrate-prime decomposition, as a backdrop to
the irrep-density observation.

We do not claim a derivation of $\Mtwentytwo$ from the substrate;
the question of whether the non-genericity reflects a structural
identification or a deep coincidence is left as the paper's central
open question. The verification is computational and runs in under
$5$ seconds with \texttt{numpy} and \texttt{sympy} as the only
external dependencies.
\end{abstract}

\maketitle

\tableofcontents

%==============================================================
\section{Introduction}
%==============================================================

\subsection{Order-factorization coincidences in the literature.}
The sporadic simple groups carry distinguished arithmetic structure.
The Mathieu group $\Mtwentytwo$ has order
\[
  |\Mtwentytwo| \;=\; 2^{7} \cdot 3^{2} \cdot 5 \cdot 7 \cdot 11
                \;=\; 443\,520,
\]
involving exactly five primes. It is the automorphism group of the
unique Steiner system $S(3, 6, 22)$ on $22$ points, with $77$ blocks
of size $6$. The combinatorial design $S(3, 6, 22)$ has a rich
parameter set --- $v = 22$ points, $b = 77$ blocks, $r = 21$
replications, $\lambda_{2} = 5$ pair-replication --- each of which
factors into small primes.

The present paper catalogs a family of order-factorization
coincidences linking $\Mtwentytwo$'s combinatorial parameters to a
specific discrete substrate $(\Z/10\Z, \sig, W)$ studied in
\cite{SandersGishFourCore}. The substrate carries the same five primes
(with the same multiplicities) as $|\Mtwentytwo|$, distributed across
elementary substrate-internal quantities.

\subsection{What this paper claims.}
The paper's central observation (Section~\ref{sec:identities}) is
quantitative: the prime-factorization profile of $\Mtwentytwo$'s
$12$ irreducible representation dimensions concentrates inside the
substrate's distinguished prime monoid at a rate that is
non-generic at the level of $p \approx 10^{-6}$ under a uniform
null model on $[1, 385]$. Section~\ref{sec:identities} states the
observation precisely (Theorem~\ref{thm:nongenericity}) and
performs the null-model computation.

A second, lighter section (Section~\ref{sec:steiner-table})
collects the standard Steiner-system parameters of $S(3, 6, 22)$
into a single table, exhibiting their factorization in the
substrate primes. We present this as a substrate-flavored
re-reading of textbook Steiner-design parameters, not as a
sequence of independent ``coincidences''; the textbook content is
attributed to Conway-Sloane~\cite{ConwaySloane1999}.

We do not claim any of the following:
\begin{itemize}\setlength\itemsep{2pt}
\item A derivation of $\Mtwentytwo$ from substrate algebra alone.
\item An action of $\Mtwentytwo$ on the substrate's $10$ elements
preserving the substrate operations.
\item A theorem of moonshine type linking $\Mtwentytwo$'s character
table to substrate dynamics.
\end{itemize}

The question of whether the non-genericity in
Theorem~\ref{thm:nongenericity} reflects a structural identification
or a deep coincidence is left as the paper's central open question
(Section~\ref{sec:open}).

\subsection{Audience.}
The paper is written for the readership of the \emph{American
Mathematical Monthly}: the Mathieu / Steiner content is laid out
elementarily (Section~\ref{sec:m22-steiner}), the substrate side is
sketched (Section~\ref{sec:substrate}), and the catalog of identities
is presented as a single section (Section~\ref{sec:identities}). No
sporadic-group background beyond standard undergraduate group theory
is assumed.

%==============================================================
\section{$\Mtwentytwo$ and the Steiner system $S(3, 6, 22)$}
\label{sec:m22-steiner}
%==============================================================

We recall the standard facts about $\Mtwentytwo$ and its Steiner
system. Standard references are Conway-Sloane
\cite{ConwaySloane1999} and Aschbacher \cite{Aschbacher1994}.

\subsection{The Mathieu group $\Mtwentytwo$.}
$\Mtwentytwo$ is one of the five sporadic Mathieu groups, of order
$443\,520 = 2^{7} \cdot 3^{2} \cdot 5 \cdot 7 \cdot 11$. It acts on a
$22$-point set $\Omega_{22}$ as the automorphism group of the Steiner
system $S(3, 6, 22)$.

\subsection{The Steiner system $S(3, 6, 22)$.}
$S(3, 6, 22)$ is the unique Steiner system on $22$ points such that
every $3$-element subset of the $22$ points lies in exactly one block
(\emph{hexad}), where each block has size $6$. The defining parameters
are:
\begin{center}
\begin{tabular}{ll}
\toprule
Parameter & Value \\
\midrule
points $v$ & $22$ \\
hexads $b$ & $77$ \\
block size $k$ & $6$ \\
$t$ (Steiner) & $3$ \\
replication $r$ (blocks per point) & $21$ \\
pair-replication $\lambda_{2}$ (blocks per pair) & $5$ \\
triple-replication $\lambda_{3}$ (blocks per triple) & $1$ \\
block stabilizer order & $5\,760 = 2^{7} \cdot 3^{2} \cdot 5$ \\
point stabilizer order & $|M_{21}| = 20\,160$ \\
\bottomrule
\end{tabular}
\end{center}

\subsection{Irreducible representations of $\Mtwentytwo$.}
$\Mtwentytwo$ has $12$ irreducible complex representations of
dimensions
\[
  \{1, 21, 45, 45, 55, 99, 154, 210, 231, 280, 280, 385\}.
\]
The natural permutation representation $\C^{22}$ on $\Omega_{22}$
decomposes as $V_{1} \oplus V_{21}$, the trivial $1$-dim and the
$21$-dim irrep (which is irreducible by double transitivity of the
$\Mtwentytwo$ action on $\Omega_{22}$).

%==============================================================
\section{The substrate}
\label{sec:substrate}
%==============================================================

\subsection*{Lens and substrate (ownership).}
The substrate $(\Z/10\Z, \sig, W)$ and its distinguished prime set
$\{2, 3, 5, 7, 11\}$ are defined in \cite{SandersGishFourCore}. The
roles of each prime are stated in the present \S below from
intrinsic substrate data --- CRT residue characteristics, $\sig^{2}$-cycle
order, HARMONY-index appearance, the threshold $T^{*} = 5/7$
numerator/denominator, and the wobble fraction's $11$-prolongation.
The non-genericity result (Theorem~\ref{thm:nongenericity}) is a
quantitative claim about $\Mtwentytwo$'s irrep-dimension factorization
profile relative to the substrate prime monoid; it does not depend
on the deeper TIG-framework interpretation of the substrate. A
reader can verify the theorem from the present \S's
intrinsic-prime definitions together with the \S\ref{sec:identities}
binomial-tail computation, without consulting \cite{SandersGishFourCore}.
The closest published precedent for the substrate's algebraic
neighborhood (small finite commutative non-associative structures) is
\cite{DrapalWanless2021}, where the opposite extremum (maximally
non-associative quasigroups) is investigated.

\subsection{The substrate primes (defined from the substrate's own data).}
The discrete substrate of \cite{SandersGishFourCore} is
$(\Z/10\Z, \sig, W)$, where $\sig$ is the permutation
$(0)(3)(8)(9)(1\,7\,6\,5\,4\,2)$ of $\Z/10\Z$ and $W = 3/50$ is the
wobble fraction. We note that $\sig$ is not an involution: it has
order $6$ on the $6$-cycle and order $1$ on the four fixed indices
$\{0, 3, 8, 9\}$, so $\sig$ has order $6$ as an element of $S_{10}$.

The substrate distinguishes the prime set
$\{2, 3, 5, 7, 11\}$ as follows; each prime is named by an
intrinsic substrate quantity, with no external input:
\begin{itemize}\setlength\itemsep{2pt}
\item $2$: the residue characteristic of the $\Z/2$ factor in the
CRT decomposition $\Z/10\Z \cong \Z/2 \times \Z/5$.
\item $3$: the order of $\sig^{2}$ restricted to the $\sig$-cycle
(direct check: $\sig^{2}$ has cycle structure $(0)(3)(8)(9)(1\,6\,4)
(7\,5\,2)$, hence order $3$ on the cycle).
\item $5$: the residue characteristic of the $\Z/5$ factor in the
CRT decomposition; equivalently, the size of the largest non-trivial
$\sig$-orbit divided by $\gcd(6, 5) = 1$.
\item $7$: the size of the unique $\sig$-orbit modulo $5$ (the
$6$-cycle has $5$ residues mod $5$ entering once and $0$ entering
once); equivalently, $|\{0, \ldots, 9\}| - |\{0, 3, 8\}| = 7$ enumerates
the non-fixed indices outside $\{8\}$. Within
\cite{SandersGishFourCore}, $7$ is also the value of the
substrate's HARMONY operator and the denominator of the threshold
$T^{*} = 5/7$.
\item $11$: the prime $11$ appears in the substrate as the
denominator of $W = 3/50$ after passing to the canonical form
$W \cdot 11/11 = 33/550$ in
\cite{SandersGishFourCore}~\S$3$ (the wobble's $11$-prolongation; the
prime $11$ is then the smallest prime not appearing in
$\{2, 3, 5, 7\}$ that the substrate's auxiliary data carries). We
note honestly that the case for $11$'s substrate-distinguishedness is
weakest of the five; the main result of this paper does not rest on
$11$ alone.
\end{itemize}

\begin{remark}
The substrate-prime distinction $\{2, 3, 5, 7, 11\}$ is a property
of the substrate's intrinsic data above. It happens that
$|\Mtwentytwo| = 2^{7} \cdot 3^{2} \cdot 5 \cdot 7 \cdot 11$ uses
exactly these five primes; this paper's main result
(Theorem~\ref{thm:nongenericity}) makes the agreement quantitative.
\end{remark}

\subsection{The wobble fraction.}
The wobble fraction is $W = 3/50$. The substrate paper
\cite{SandersGishFourCore} also defines a complementary fraction
$G = 22/50$ ($G$ for the gentleness phase amplitude), with
$W + G = 25/50 = 1/2$. This identity of fractions is internal to
the substrate; the gentleness numerator $22$ and the size of the
Mathieu action set $|\Omega_{22}| = 22$ are equal as integers but not
related by any structural map in the present paper.

%==============================================================
\section{The non-generic concentration of irrep dimensions}
\label{sec:identities}
%==============================================================

We state and prove the paper's main observation: the prime
factorizations of $\Mtwentytwo$'s irrep dimensions concentrate
inside the substrate's prime monoid at a rate that is non-generic
under a uniform null on $[1, 385]$.

\begin{definition}[Substrate-prime band]
\label{def:band}
For an integer $m \ge 1$, write $m \in \mathcal{B}$ (the
\emph{substrate-prime band}) iff every prime factor of $m$ lies in
$\{2, 3, 5, 7, 11\}$ and $\nu_{2}(m) \le 1$ (i.e., $m$ has at most a
single factor of $2$).
\end{definition}

We classify $\Mtwentytwo$'s $12$ irrep dimensions by membership in
the band:

\begin{center}
\begin{tabular}{r l c c}
\toprule
$\dim V$ & factorization & strict $\{3,5,7,11\}$? & $\in \mathcal{B}$? \\
\midrule
$1$    & trivial               & vacuously yes & yes \\
$21$   & $3 \cdot 7$           & yes & yes \\
$45$   & $3^{2} \cdot 5$       & yes & yes \\
$45$   & $3^{2} \cdot 5$       & yes & yes \\
$55$   & $5 \cdot 11$          & yes & yes \\
$99$   & $3^{2} \cdot 11$      & yes & yes \\
$154$  & $2 \cdot 7 \cdot 11$  & no (has $2$) & yes \\
$210$  & $2 \cdot 3 \cdot 5 \cdot 7$ & no & yes \\
$231$  & $3 \cdot 7 \cdot 11$  & yes & yes \\
$280$  & $2^{3} \cdot 5 \cdot 7$ & no & no ($2^{3}$) \\
$280$  & $2^{3} \cdot 5 \cdot 7$ & no & no ($2^{3}$) \\
$385$  & $5 \cdot 7 \cdot 11$  & yes & yes \\
\bottomrule
\end{tabular}
\end{center}

The factorizations are direct calculations (verified in
\texttt{m22\_decomposition.py}, \S\ref{sec:verification}). The
dimensions are taken from the standard $\Mtwentytwo$ character table
(Conway--Sloane~\cite{ConwaySloane1999}; ATLAS~\cite{ATLAS1985}).
Sum-of-squares verification: $\sum (\dim V_{i})^{2} = 443\,520 =
|\Mtwentytwo|$ (checked directly).

The classification yields:
\begin{itemize}\setlength\itemsep{2pt}
\item Strict $\{3, 5, 7, 11\}$ (no factor of $2$):
\textbf{eight} dimensions, of which seven are non-trivial:
$21$, $45$ (with multiplicity $2$), $55$, $99$, $231$, $385$.
\item Dimensions in $\mathcal{B}$ (the substrate-prime band with
at most one factor of $2$): \textbf{ten} of $12$, including the
trivial $1$, the seven strict non-trivial dimensions, and
$154 = 2 \cdot 7 \cdot 11$ and $210 = 2 \cdot 3 \cdot 5 \cdot 7$.
\item Dimensions outside $\mathcal{B}$: only the two copies of
$280 = 2^{3} \cdot 5 \cdot 7$.
\end{itemize}

\begin{theorem}[Non-genericity of $\Mtwentytwo$ irrep concentration]
\label{thm:nongenericity}
Let $\mathcal{B}_{N}$ denote the integers in $[1, N]$ that lie in
the substrate-prime band $\mathcal{B}$ (Definition~\ref{def:band}).
Then:
\begin{enumerate}\setlength\itemsep{4pt}
\item $|\mathcal{B}_{385}| = 67$, giving null density
$|\mathcal{B}_{385}| / 385 \approx 0.1740$.
\item $\Mtwentytwo$ has $10$ of $12$ irrep dimensions in
$\mathcal{B}$, observed density $10/12 \approx 0.8333$.
\item Under the null hypothesis $H_{0}$ that the $12$ irrep
dimensions are i.i.d.\ uniform on $\{1, 2, \ldots, 385\}$,
\[
  P[X \ge 10 \mid X \sim \mathrm{Bin}(12, 0.1740)]
  \;\approx\; 1.19 \times 10^{-6}.
\]
\item Under the same null, the $p$-value for the strict
$\{3, 5, 7, 11\}$ count of seven non-trivial dimensions (out of $11$
non-trivial draws) is
$P[X \ge 7 \mid X \sim \mathrm{Bin}(11, 0.0990)] \approx 2.14
\times 10^{-5}$.
\end{enumerate}
\end{theorem}

\begin{proof}
\textit{(1)} Direct enumeration:
$\mathcal{B}_{385} = \{1, 2, 3, 5, 6, 7, 9, 10, 11, 14, 15, 18, 21,
22, 25, 27, 30, 33, 35, 42, 45, 49, 50, 54, 55, 63, 66, 70, 75, 77,
\ldots, 385\}$ has $67$ elements; verified by the one-liner
\[
  \texttt{[n for n in range(1,386) if all(p in \{2,3,5,7,11\}
  for p in factorint(n)) and factorint(n).get(2,0) <= 1]}
\]
in \texttt{m22\_decomposition.py}.

\textit{(2)} The classification table above; verified by direct
factorization of each dimension.

\textit{(3)} The binomial tail probability:
\[
  P[X \ge 10] \;=\; \sum_{k=10}^{12} \binom{12}{k}
  (0.1740)^{k} (0.8260)^{12-k}
  \;\approx\; 1.19 \times 10^{-6}.
\]
The numerical value is computed in \texttt{m22\_decomposition.py}.

\textit{(4)} Among integers in $[2, 385]$, the count factoring strictly
in $\{3, 5, 7, 11\}$ (no factor of $2$) is $38$, density $38 / 384
\approx 0.0990$. The seven non-trivial irrep dimensions in this strict
band (out of $11$ non-trivial irreps) yield
$P[X \ge 7 \mid X \sim \mathrm{Bin}(11, 0.0990)] \approx 2.14 \times
10^{-5}$.
\end{proof}

\begin{remark}[Status of the result]
The non-genericity of Theorem~\ref{thm:nongenericity} is a
quantitative statement, not a structural identification. The null
model is the simplest reasonable choice (uniform on $[1, 385]$);
alternative null models (e.g., conditioning on a sum-of-squares
constraint, or restricting to $385$-divisors) would yield different
$p$-values. The point of the theorem is that under the natural
null, the concentration is real and small.

A genuine derivation of $\Mtwentytwo$'s irrep dimensions from
substrate-internal data would close the question; the present
result does not attempt this.
\end{remark}

\begin{remark}[The dimension $231$]
Among the $12$ $\Mtwentytwo$-irrep dimensions, $231 = 3 \cdot 7
\cdot 11$ is the unique one whose prime factorization is a
squarefree product of exactly three distinct substrate primes
(no $2$, no $5$). This factorization-pattern uniqueness identifies
$231$ as the most substrate-prime-saturated dimension; it sharpens
Theorem~\ref{thm:nongenericity} but does not extend it.
\end{remark}

%==============================================================
\section{Steiner-system parameters of $S(3, 6, 22)$ as backdrop}
\label{sec:steiner-table}
%==============================================================

The standard parameters of the Steiner system $S(3, 6, 22)$
(Conway--Sloane~\cite{ConwaySloane1999};
Cameron~\cite{Cameron1999}) factor in small primes; the
substrate-prime decomposition is recorded in the rightmost column.
We present this as a substrate-flavored re-reading of textbook
quantities, not as independent ``coincidences'' alongside the main
result of \S\ref{sec:identities}.

\begin{center}
\begin{tabular}{l c c l}
\toprule
Parameter & Symbol & Value & Substrate-prime decomposition \\
\midrule
points     & $v$         & $22$  & $2 \cdot 11$ \\
hexads     & $b$         & $77$  & $7 \cdot 11$ (HARMONY $\times$ wobble-prime) \\
block size & $k$         & $6$   & $2 \cdot 3$ \\
Steiner $t$ & $t$        & $3$   & $3$ (substrate trinitarian factor) \\
replication & $r$        & $21$  & $3 \cdot 7$ ($= \dim V_{21}$) \\
pair-replication & $\lambda_{2}$ & $5$ & $5$ ($= T^{*}$ numerator) \\
triple-replication & $\lambda_{3}$ & $1$ & trivial \\
block stabilizer order & --- & $5\,760$ & $2^{7} \cdot 3^{2} \cdot 5$ \\
point stabilizer order & $|M_{21}|$ & $20\,160$ & $2^{6} \cdot 3^{2} \cdot 5 \cdot 7$ \\
\bottomrule
\end{tabular}
\end{center}

The factorizations are textbook (Conway--Sloane~\cite{ConwaySloane1999}
Ch.~$10$--$11$; ATLAS~\cite{ATLAS1985}; Wilson~\cite{Wilson2009}). The
substrate-prime decomposition is direct.

\begin{remark}[The substrate's $\sig$-orbit and the hexad block size]
The substrate $\sig$-orbit on $\{1, 7, 6, 5, 4, 2\}$ has size $6$,
the block size $k$ of $S(3, 6, 22)$. The integers match. We do not
conclude that the $\sig$-orbit is a distinguished hexad in
$S(3, 6, 22)$; the Steiner system is unique up to isomorphism, and
selecting a specific isomorphism class with the $\sig$-orbit as a
distinguished hexad would be a substantive representation-theoretic
question (Section~\ref{sec:open}).
\end{remark}

%==============================================================
\section{Computational verification}
\label{sec:verification}
%==============================================================

All numerical claims of \S\ref{sec:identities} and
\S\ref{sec:steiner-table} are verified computationally:

\begin{itemize}\setlength\itemsep{2pt}
\item \texttt{m22\_decomposition.py} factors each
$\Mtwentytwo$-irrep dimension, classifies each dimension by
membership in the substrate-prime band $\mathcal{B}$, enumerates
$\mathcal{B}_{385}$ to obtain the null density, and computes the
binomial tail probabilities of
Theorem~\ref{thm:nongenericity}(3)--(4) using \texttt{math.comb}.
\item \texttt{steiner\_sigma\_hexad.py} confirms the Steiner-system
parameters $(v, b, k, r, \lambda_{2}, \lambda_{3}) =
(22, 77, 6, 21, 5, 1)$ via the formulas $b = \binom{v}{t}/\binom{k}{t}$,
$r = bk/v$, $\lambda_{i} = \binom{k-i}{t-i} r / \binom{v-i}{t-i}$.
\end{itemize}

All scripts run in under $5$ seconds with \texttt{sympy} (and
optionally \texttt{numpy}) as the only external dependencies.
Scripts deposited at the public repository
\url{https://github.com/TiredofSleep/ck/tree/tig-synthesis} (DOI
\texttt{10.5281/zenodo.18852047}).

%==============================================================
\section{Open questions}
\label{sec:open}
%==============================================================

\begin{enumerate}
\item \textbf{Structural origin.} Is the non-genericity of
Theorem~\ref{thm:nongenericity} the shadow of a structural
identification of $\Mtwentytwo$ with substrate-internal data, or is
it a deep arithmetic coincidence with no algebraic content? The
paper does not commit to either reading.

\item \textbf{$\Mtwentytwo$ action on the substrate.} Does
$\Mtwentytwo$ act on the substrate $(\Z/10\Z, \sig)$ in a way
respecting the magma operations $\TS, \BH$ of
\cite{SandersGishFourCore}? The substrate has only $10$ elements,
far fewer than the $22$ on which $\Mtwentytwo$ acts; any action
would necessarily involve the substrate's wobble layer or some
auxiliary structure.

\item \textbf{Structural Steiner-class selection.} $S(3, 6, 22)$ is
unique up to isomorphism, but the embedding of the $\sig$-orbit as
a specific hexad is a labeling choice. Does the substrate algebra
of \cite{SandersGishFourCore} select a canonical representative
isomorphism class with the $\sig$-orbit as a distinguished hexad?

\item \textbf{Relation to Mathieu moonshine.} The Mathieu groups
appear in moonshine phenomena (notably $M_{24}$ in K3-surface
moonshine; see \cite{Eguchi2011}). The present non-genericity is
arithmetic, not analytic (no generating function). Does it connect
to any moonshine context?

\item \textbf{Other sporadic groups.} The substrate primes
$\{2, 3, 5, 7, 11\}$ are also a subset of the prime divisors of
$|M_{11}|$, $|M_{12}|$, $|M_{23}|$, $|M_{24}|$. Does the
non-genericity extend in a structurally illuminating way to the
larger Mathieu groups?
\item \textbf{Tighter null model.} Theorem~\ref{thm:nongenericity}
uses the simplest null (uniform on $[1, 385]$). A tighter null
conditioning on the sum-of-squares constraint
$\sum (\dim V_{i})^{2} = 443\,520$ would refine the $p$-value; the
authors expect the qualitative non-genericity to persist but have
not computed it.
\end{enumerate}

%==============================================================
\section*{Acknowledgments}
%==============================================================

This work catalogs order-factorization identities between the Mathieu
group $\Mtwentytwo$ and a discrete substrate of independent algebraic
interest. The identities are presented as a self-contained
arithmetic-combinatorial puzzle for the \emph{Monthly} readership.
The verification scripts and substrate-side definitions are part of
the broader Trinity Infinity Geometry framework developed at 7Site
LLC (public repository at github.com/TiredofSleep/ck, DOI
10.5281/zenodo.18852047).

%==============================================================
\begin{thebibliography}{99}
%==============================================================

\bibitem{ConwaySloane1999} J.~H.~Conway, N.~J.~A.~Sloane,
\emph{Sphere Packings, Lattices and Groups} (3rd ed.),
Grundlehren der mathematischen Wissenschaften 290, Springer, 1999.
[Standard reference for the Mathieu groups and the Steiner systems
$S(5, 6, 12)$, $S(5, 8, 24)$.]

\bibitem{ATLAS1985} J.~H.~Conway, R.~T.~Curtis, S.~P.~Norton,
R.~A.~Parker, R.~A.~Wilson,
\emph{ATLAS of Finite Groups: Maximal Subgroups and Ordinary
Characters for Simple Groups},
Oxford University Press, 1985.
[Reference for $\Mtwentytwo$ character table and irrep dimensions.]

\bibitem{Aschbacher1994} M.~Aschbacher,
\emph{Sporadic Groups},
Cambridge Tracts in Mathematics 104, Cambridge University Press, 1994.

\bibitem{Cameron1999} P.~J.~Cameron,
\emph{Permutation Groups},
London Mathematical Society Student Texts 45, Cambridge University
Press, 1999.

\bibitem{Wilson2009} R.~A.~Wilson,
\emph{The Finite Simple Groups},
Graduate Texts in Mathematics 251, Springer, 2009.

\bibitem{Diaconis1988} P.~Diaconis,
\emph{Group Representations in Probability and Statistics},
IMS Lecture Notes -- Monograph Series 11, Institute of Mathematical
Statistics, 1988.
[Canonical \emph{Monthly}-level exposition of finite-Fourier
arithmetic and group representations.]

\bibitem{Eguchi2011} T.~Eguchi, H.~Ooguri, Y.~Tachikawa,
\emph{Notes on the K3 surface and the Mathieu group $M_{24}$},
Exp.~Math.~\textbf{20} (2011), 91--96.

\bibitem{DrapalWanless2021} A.~Dr{\'a}pal, I.~M.~Wanless,
\emph{Maximally nonassociative quasigroups},
J.~Combin.\ Theory Ser.~A \textbf{184} (2021), 105510.
[Domain precedent: small finite commutative non-associative
structures; opposite extremum to the substrate of
\cite{SandersGishFourCore}.]

\bibitem{SandersGishFourCore} B.~R.~Sanders, M.~Gish,
\emph{Joint Closure, Per-Coordinate Fuse Data, and a Closed-Form
Algebraic Attractor of Two Commutative Binary Operations on
$\Z/10\Z$}, submitted to \emph{Algebraic Combinatorics}, 2026 (J02).

\bibitem{TIGsynthesis2026} B.~R.~Sanders,
\emph{Trinity Infinity Geometry Synthesis},
github.com/TiredofSleep/ck (default branch),
DOI: 10.5281/zenodo.18852047, 2026.

\end{thebibliography}

\end{document}
